% AQIDA_CONSTITUTIONAL_SCOPE: gate
% THEORY_BINDING_MAP:
% - CORE_LOGIC v3, A9 and A10 (an energy has no cross-cancellation): Theorem 1 uses a positive form exactly
%   once (an isometry does not increase norms; for vectors: Cauchy-Schwarz). Section 5 shows that the complex
%   bilinear (Minkowski) representation satisfies every other step. Section 6 uses energies tr(X* g X) >= 0.
% - CORE_LOGIC v3, section 2.2 (a zero carries its witness): the space N of Section 3 consists of the
%   coefficient arrays whose restriction to every layer is a linear relation of that layer.
% - n_s_u: explicit mechanism rows: Table 1.  elm_eulerian, e_i/squaring/AQiDA_Physics/Chess_Cybertron: not used.
% CLAIM_BOUNDARY: unified manuscript of gate row 19 of DRIFT_LOCK.md. Claims: Theorem 1 and its corollaries
%   (a bound of order t^2 for the real, Hermitian and projective dimension of the triple graph, which is what
%   Problem 1 of ECCC TR26-151 asks for), and the linear bounds of Section 6. Verification status is stated in
%   Section 7 and in the footnote of the title; nothing here has been peer reviewed. No claim about Hurwitz's
%   problem.
% FALSIFICATION_TESTS: the Lean development of Section 7; cube_layer_theorem_check.py and
%   cube_pointwise_bound_check.py (the Minkowski Gram matrix must violate the bound); the certificate files.
\documentclass[11pt]{article}
\usepackage[a4paper,margin=1in]{geometry}
\usepackage[T1]{fontenc}
\usepackage[utf8]{inputenc}
\usepackage{lmodern}
\usepackage{amsmath,amssymb,amsthm,mathtools}
\usepackage{booktabs}
\usepackage[colorlinks=true,linkcolor=blue,citecolor=blue,urlcolor=blue]{hyperref}

\newtheorem{theorem}{Theorem}
\newtheorem{corollary}[theorem]{Corollary}
\newtheorem{proposition}[theorem]{Proposition}
\newtheorem{lemma}[theorem]{Lemma}
\theoremstyle{remark}
\newtheorem{remark}[theorem]{Remark}

\newcommand{\R}{\mathbb{R}}
\newcommand{\C}{\mathbb{C}}
\newcommand{\K}{\mathbb{K}}
\newcommand{\ip}[2]{\langle #1,#2\rangle}
\newcommand{\norm}[1]{\lVert #1\rVert}
\newcommand{\rank}{\operatorname{rank}}
\newcommand{\tr}{\operatorname{tr}}
\newcommand{\floor}[1]{\lfloor #1\rfloor}
\newcommand{\Gt}{G_t}
\newcommand{\Hn}{H_n}

\title{A quadratic lower bound for real orthogonal representations\\ of the graph of triples\thanks{Draft of
26 November 2026. Produced with an AI system (see the statement before the references). Not peer reviewed.
Verification status: Section~\ref{sec:lean}.}}
\author{Babar Ahmed Khan}
\date{}

\begin{document}
\maketitle

\begin{abstract}
Let $\Gt$ be the graph whose vertices are the $3$-subsets of $[t]$, two of them being adjacent when they share
exactly one element (the generalised Johnson graph $J(t,3,1)$). An orthogonal representation assigns a non-zero
vector to each vertex so that adjacent vertices get orthogonal vectors. Hrube\v{s} and Iyer asked for a proof
that over the real numbers the least dimension $d_\R(t)$ is of order $t^2$; over $\C$ with the bilinear form
dimension $t+1$ suffices, and the Lov\'asz number of the graph, the classical lower bound, is about $3t/4$.
We prove
\[ d_\R(t)\;\ge\;\tfrac15\,\floor{t/3}^2 , \]
so that $d_\R(t)=\Theta(t^2)$. The proof is one page of linear algebra: it counts linear relations inside the
layers of a three-dimensional cube and uses the positivity of the inner product once. The same bound holds for
complex vectors with the Hermitian form and for representations by projectors of any rank, hence for the
projective rank and the quantum chromatic number of $J(t,3,1)$. The complex bilinear representation satisfies
every step of the proof except the one that uses positivity. For small $t$, where the quadratic bound is weak,
we give bounds of a different kind: an explicit bound $(\sqrt3-o(1))\,t$, the inequality $d_\R(t)\ge t+2$ for
every $t\ge 11$, and exact certificates for $t=10,\dots,16$. The quadratic bounds (for real and Hermitian
vectors and for subspaces) and the upper bound $t^2$ have been formalised in Lean~4 with Mathlib; the linear
bounds have not.
\end{abstract}

\section{Introduction}

\subsection{The problem}
For a field $\K\in\{\R,\C\}$ a \emph{representation} of $\Gt$ in dimension $d$ is a family of non-zero vectors
$v_A\in\K^d$, one for each $3$-subset $A$ of $[t]$, with $\ip{v_A}{v_B}=0$ whenever $|A\cap B|=1$; here
$\ip{\cdot}{\cdot}$ is the standard inner product (Hermitian over $\C$). We write $d_\R(t)$ and $d_{\mathrm H}(t)$
for the least dimension over $\R$ and over $\C$. In the notation of Hrube\v{s} and Iyer~\cite{HI},
$d_\R(t)=\dim_\R(\overline{J(t,3,1)})$. Their Problem~1 reads: ``Show that
$\dim_\R(\overline{J(n,3,1)})\ge\Omega(n^2)$'' (their $n$ is our $t$). They record the trivial bounds
$\Omega(t)\le\dim_\R\le O(t^2)$, the bound $t+1$ for the bilinear form over $\C$ (realised by the vectors
$(1_A,i)$, for which $(1_A,1_B)-1=|A\cap B|-1$), and a bound $\Omega(t^2)$ over fields of characteristic two.
Their method for real representations of Cayley graphs is a linear programme whose value is the Lov\'asz
number~\cite{Lo} \cite[Lemma~22 and~(8)]{HI}. For $\Gt$ the Lov\'asz bound is
$1+3\binom{t-3}2/(2t-11)\approx 3t/4$ for $t\ge14$ (our computation). The upper bound $O(t^2)$ is elementary: give each triple the basis vector indexed by one
of its pairs.

\subsection{Results}
The \emph{cube graph} $\Hn$ has vertex set $[n]^3$; two points are adjacent when they agree in exactly one
coordinate. If $[t]$ contains three disjoint $n$-sets $X,Y,Z$, the triples $\{x,y,z\}$ with $x\in X$, $y\in Y$,
$z\in Z$ form a copy of $\Hn$ inside $\Gt$: two of them share exactly one element exactly when they agree in
exactly one coordinate. A \emph{layer} of the cube is a set $\{p: p_i=a\}$ ($3n$ layers, $n^2$ points each).

\begin{theorem}\label{thm:L}
Let $\K\in\{\R,\C\}$, let $E$ be a finite-dimensional inner-product space over $\K$ and $r\ge1$. Suppose that
to every point $p\in[n]^3$ there is assigned a linear isometry $U_p\colon\K^r\to E$ such that the ranges of
$U_p$ and $U_q$ are orthogonal whenever $p$ and $q$ are adjacent in $\Hn$. For a layer $\Lambda$ let
$\rank(\Lambda)$ be the dimension of the span of the ranges of the $U_p$, $p\in\Lambda$. Then
\[ \sum_{\Lambda}\rank(\Lambda)\;\ge\;\tfrac35\,r\,n^3 ,\qquad\text{and hence}\qquad
   \dim E\;\ge\;\tfrac15\,r\,n^2 . \]
\end{theorem}

For $r=1$ this is a statement about unit vectors $u_p=U_p(1)$: \emph{if unit vectors $u_p$, $p\in[n]^3$, are
orthogonal whenever two points agree in exactly one coordinate, they span at least $n^2/5$ dimensions.}

\begin{corollary}\label{cor:triples}
For every $t$: $d_\R(t)\ge\floor{t/3}^2/5$ and $d_{\mathrm H}(t)\ge\floor{t/3}^2/5$. Consequently
$d_\R(t)=\Theta(t^2)$, which answers Problem~1 of~\cite{HI}.
\end{corollary}

\begin{corollary}\label{cor:proj}
If orthogonal projectors $P_A$ of rank $r$ on $\K^D$ are assigned to the $3$-subsets of $[t]$ with $P_AP_B=0$
whenever $|A\cap B|=1$, then $D/r\ge\floor{t/3}^2/5$. Hence the projective rank $\xi_f$ and the quantum
chromatic number $\chi_q$ satisfy
$\chi_q(J(t,3,1))\ge\xi_f(J(t,3,1))\ge\floor{t/3}^2/5$.
\end{corollary}

So the real, the Hermitian and the projective dimension of this graph are all of order $t^2$, while its
Lov\'asz number and its complex bilinear dimension are of order $t$.

\subsection{Where positivity enters}
The proof of Theorem~\ref{thm:L} uses the positivity of the form in one line: an isometry $U_p$ satisfies
$\norm{U_p^*a}\le\norm a$ (for $r=1$ this is the Cauchy--Schwarz inequality). Section~\ref{sec:mink} shows
that the complex bilinear representation, which lives in a space with one negative direction, satisfies every
other step, and violates exactly this one as soon as $n\ge4$.

\subsection{Small $t$}
Corollary~\ref{cor:triples} exceeds the complex bilinear bound $t+1$ only for $t\ge48$. For small $t$ we use a
different method (moments of the representation up to order three), which gives:
\begin{itemize}
\item an explicit bound $d\ge(\sqrt3-o(1))\,t$ valid for every $t\ge7$ (Theorem~\ref{thm:A});
\item $d_\R(t)\ge t+2$ and $d_{\mathrm H}(t)\ge t+2$ for \emph{every} $t\ge11$ (Theorem~\ref{thm:tplus2});
\item $d\ge 11,14,16,17,19,21,23$ for $t=10,\dots,16$, by certificates checked in exact arithmetic
(Theorem~\ref{thm:cert}).
\end{itemize}
These bounds hold for projectors of every rank, and Proposition~\ref{prop:proj} shows that for that reason the
method cannot give more than $25/2$ at $t=10$. Table~\ref{tab:bounds} compares the bounds; the closed bound of
Theorem~\ref{thm:A} is the better one up to $t=74$, Corollary~\ref{cor:triples} from $t=75$ on.

\begin{table}[ht]\centering
\begin{tabular}{rrrrrrl}
\toprule
$t$ & Lov\'asz & Thm.~\ref{thm:A} & Thm.~\ref{thm:cert} & Cor.~\ref{cor:triples} & $t+1$ & best representation known\\
\midrule
9 & 7.43 & 7.99 & -- & 1.8 & 10 & 8 (exact value)\\
10 & 8.00 & 9.56 & 11 & 1.8 & 11 & 15\\
11 & 8.64 & 11.16 & 14 & 1.8 & 12 & 18 (numerical)\\
12 & 9.31 & 12.79 & 16 & 3.2 & 13 & \\
13 & 10.00 & 14.44 & 17 & 3.2 & 14 & \\
16 & 12.14 & 19.47 & 23 & 5.0 & 17 & \\
48 & 35.94 & 74.60 & -- & 51.2 & 49 & \\
78 & 58.41 & 126.55 & -- & 135.2 & 79 & \\
100 & 74.90 & 164.65 & -- & 217.8 & 101 & \\
1000 & 749.88 & 1723.51 & -- & 22177.8 & 1001 & at most $\binom{1000}{2}$\\
\bottomrule
\end{tabular}
\caption{Columns 2 to 5: lower bounds for the dimension of a representation of $\Gt$ by real or Hermitian unit
vectors. Column 6: the dimension of the complex bilinear representation. Column 7: the smallest dimension of a
real representation that we know.}
\label{tab:bounds}
\end{table}

\section{Notation}
$P=[n]^3$. For $p\in P$ and $i\in\{1,2,3\}$, the \emph{line} $\ell_i(p)$ is the set of the $n$ points that agree
with $p$ outside coordinate $i$. Two distinct points of a line agree in two coordinates; a point $q\ne p$ of
$\ell_i(p)$ and a point $q'\ne p$ of $\ell_j(p)$, $i\ne j$, agree in exactly one coordinate. The layer through
$p$ in which coordinate $h$ is fixed contains $\ell_i(p)\cup\ell_j(p)$ ($\{i,j,h\}=\{1,2,3\}$), and each of its
other points agrees with $p$ in coordinate $h$ only. Inner products are linear in the first argument. An
isometry $U\colon\K^r\to E$ satisfies $U^*U=I_r$, and the ranges of $U_p$, $U_q$ are orthogonal exactly when
$U_p^*U_q=0$.

\section{Proof of Theorem~\ref{thm:L}}\label{sec:proof}
An \emph{array} is an element $v=(v_p)_{p\in P}$ of $(\K^r)^P$, with $\norm v^2=\sum_p\norm{v_p}^2$. For a set
$S\subseteq P$ put $\sigma_S(v)=\sum_{p\in S}U_pv_p\in E$.

\smallskip\noindent\textbf{Step 1 (relations).} Let $N$ be the space of arrays $v$ with $\sigma_\Lambda(v)=0$
for every layer $\Lambda$. It is the kernel of a linear map from $(\K^r)^P$ into the direct sum, over the
layers, of the spans of the layers, so
\begin{equation}\label{eq:rn}
\dim N\;\ge\; r\,n^3-\sum_\Lambda\rank(\Lambda).
\end{equation}

\noindent\textbf{Step 2 (the value $\tfrac12$).} Let $v\in N$, $p\in P$, and $a_i=\sigma_{\ell_i(p)}(v)$ for
$i=1,2,3$. Then
\begin{equation}\label{eq:half}
U_p^*a_i=\tfrac12\,v_p\qquad(i=1,2,3).
\end{equation}
Indeed, let $\Lambda$ be the layer through $p$ with coordinate $3$ fixed. Applying $U_p^*$ to
$\sigma_\Lambda(v)=0$ gives $\sum_{q\in\Lambda}U_p^*U_qv_q=0$. The points of $\Lambda$ outside
$\ell_1(p)\cup\ell_2(p)$ are adjacent to $p$ and contribute nothing; $p$ lies on both lines and
$U_p^*U_p=I_r$. So $U_p^*a_1+U_p^*a_2=v_p$. The two other layers through $p$ give
$U_p^*a_1+U_p^*a_3=v_p$ and $U_p^*a_2+U_p^*a_3=v_p$, and the three equations give~\eqref{eq:half}.

\smallskip\noindent\textbf{Step 3 (norms and orthogonality).} For every line $\ell$ and $a=\sigma_\ell(v)$,
applying~\eqref{eq:half} at each point $q$ of $\ell$ (with the same line),
\begin{equation}\label{eq:norm}
\norm a^2=\sum_{q\in\ell}\ip{a}{U_qv_q}=\sum_{q\in\ell}\ip{U_q^*a}{v_q}=\tfrac12\sum_{q\in\ell}\norm{v_q}^2 .
\end{equation}
For $i\ne j$ we have $\ip{a_i}{a_j}=0$: in
$\ip{a_i}{a_j}=\sum_{q\in\ell_i(p)}\sum_{q'\in\ell_j(p)}\ip{U_qv_q}{U_{q'}v_{q'}}$ the terms with $q\ne p\ne q'$
vanish, because such $q,q'$ are adjacent; the terms with $q=p$ add up to
$\ip{U_pv_p}{a_j}=\ip{v_p}{U_p^*a_j}=\tfrac12\norm{v_p}^2$, those with $q'=p$ to
$\ip{a_i}{U_pv_p}=\tfrac12\norm{v_p}^2$, and the term $q=q'=p$, which lies in both groups and is therefore subtracted once, is $\norm{v_p}^2$; the
total is $(\tfrac12+\tfrac12-1)\norm{v_p}^2=0$.

\smallskip\noindent\textbf{Step 4 (positivity).} Let $a=a_1+a_2+a_3$. By~\eqref{eq:half},
$U_p^*a=\tfrac32v_p$. By Step~3 and since the three lines through $p$ meet only in $p$,
\[ \norm a^2=\tfrac12\sum_{i=1}^3\sum_{q\in\ell_i(p)}\norm{v_q}^2\;\le\;\tfrac12\bigl(\norm v^2+2\norm{v_p}^2\bigr). \]
Because $U_p$ is an isometry, $\norm{U_p^*a}\le\norm a$. Hence
$\tfrac94\norm{v_p}^2\le\tfrac12\norm v^2+\norm{v_p}^2$, that is,
\begin{equation}\label{eq:pb}
5\,\norm{v_p}^2\;\le\;2\,\norm v^2\qquad\text{for every }v\in N\text{ and every }p\in P .
\end{equation}

\noindent\textbf{Step 5 (dimension).} Let $\Pi$ be the orthogonal projection of $(\K^r)^P$ onto $N$ and
$\Pi_p$ its diagonal block at $p$, a positive semidefinite $r\times r$ matrix. For $c\in\K^r$ the array
$v=\Pi(c\text{ at }p)$ lies in $N$, with $v_p=\Pi_pc$ and $\norm v^2=\ip{\Pi_pc}{c}$. By~\eqref{eq:pb},
$5\norm{\Pi_pc}^2\le2\ip{\Pi_pc}{c}$, so every eigenvalue of $\Pi_p$ is at most $\tfrac25$ and
$\tr\Pi_p\le\tfrac25r$. Therefore $\dim N=\tr\Pi=\sum_p\tr\Pi_p\le\tfrac25\,r\,n^3$.

\smallskip\noindent\textbf{Step 6.} With~\eqref{eq:rn}: $\sum_\Lambda\rank(\Lambda)\ge\tfrac35rn^3$. There are
$3n$ layers, each of rank at most $\dim E$; so $\dim E\ge\tfrac15rn^2$. \qed

\begin{remark}
For the $27$ vectors of Section~\ref{sec:mink} at $n=3$ and $v$ the projection of a basis array, the
inequality $\norm{U_p^*a}\le\norm a$ of Step~4 is an equality; the loss is in the estimate of $\norm a^2$, and
$\tr\Pi_p=8/27<\tfrac25$ there. A first version of the argument, which applied Pythagoras to the three
subspaces $\{\sigma_{\ell_i(p)}(v):v\in N\}$, gave the constant $1/9$ in place of $1/5$. We did not try to
optimise the constant; for $\Hn$ the truth lies between $n^2/5$ and $n^2$ (one colour for each line of a fixed
direction), and between $n^2/5$ and $2n^2/3$ when $3$ divides $n$.
\end{remark}

\section{Consequences}

\begin{proof}[Proof of Corollary~\ref{cor:triples}]
Let $n=\floor{t/3}$ and choose disjoint $n$-sets $X,Y,Z\subseteq[t]$. Normalise the vectors of the triples
$\{x,y,z\}$, $x\in X$, $y\in Y$, $z\in Z$; they are unit vectors indexed by $[n]^3$, orthogonal whenever two
points agree in exactly one coordinate. Theorem~\ref{thm:L} with $r=1$ gives $d\ge n^2/5$.
\end{proof}

\begin{proof}[Proof of Corollary~\ref{cor:proj}]
A projector of rank $r$ is $U U^*$ for an isometry $U\colon\K^r\to\K^D$, and $P_AP_B=0$ means that the ranges
are orthogonal. Restricting to the transversal triples as above, Theorem~\ref{thm:L} gives $D\ge rn^2/5$.
The projective rank $\xi_f(G)$ is the infimum of $D/r$ over such assignments of projectors to the vertices of
$G$ with orthogonal projectors on adjacent vertices, and $\xi_f(G)\le\chi_q(G)$; see Section~\ref{sec:lit}
for the sources and conventions.
\end{proof}

\begin{remark}[More coordinates]\label{rem:k}
The same proof applies to $[n]^k$, $k\ge3$, with orthogonality required for points at Hamming distance exactly
two, taking as layers the $\binom k2n^{k-2}$ two-dimensional planes. The $k$ line sums satisfy
$U_p^*a_i=\tfrac12v_p$, Step~4 gives $(k^2-2k+2)\norm{v_p}^2\le2\norm v^2$, and one obtains
$\dim E\ge\frac{2(k-2)}{(k-1)(k^2-2k+2)}\,r\,n^2$. For the graph $J(t,k,k-2)$ this is a bound of order $t^2$
for each fixed $k$.
\end{remark}

\section{The complex bilinear representation}\label{sec:mink}
On $[n]^3$ consider the symmetric bilinear form with Gram matrix
$G(p,q)=\tfrac12(\#\{i:p_i=q_i\}-1)$. It has unit diagonal and vanishes on adjacent pairs; it is the
restriction to the transversal triples of the form $(1_A,1_B)-1$ behind the bilinear representation of
dimension $t+1$. Its rank is $3n-2$ for $n\ne3$, and it has exactly one negative eigenvalue when $n>3$; for
$n=3$ it is positive semidefinite of rank $6$ (the $27$ vectors are roots of $E_6$, and for $t=9$ the $84$ vectors $1_A-\tfrac13\mathbf 1$ are roots of
$E_8$).

Define $N$ as the space of arrays whose restriction to every layer lies in the kernel of the Gram matrix of
that layer. Steps~1, 2 and~3 are identities between entries of the Gram matrix and hold verbatim. Here
$N$ consists of the arrays whose sum along every line is zero, $\dim N=(n-1)^3$, and the diagonal entries of
the projection onto $N$ equal $((n-1)/n)^3$. This exceeds $\tfrac25$ exactly when $n\ge4$, so the conclusion~\eqref{eq:pb} of
Step~4 fails there; the form is indefinite for $n\ge4$, and $\norm{U_p^*a}\le\norm a$ is the only step of the
proof that uses positivity. The proof thus separates the bilinear representation from a real one through
that single inequality.
(Numerically: $0.296$ at $n=3$; $0.422,\ 0.512,\ 0.579,\ 0.670$ at $n=4,5,6,8$.)

\section{Linear bounds for small $t$}\label{sec:linear}
The bounds of this section hold for projective representations: orthogonal projectors $P_A$ of rank $r$ on
$\R^D$ or $\C^D$ with $P_AP_B=0$ on adjacent pairs; the \emph{cost} is $d=D/r$, and vectors are the case
$r=1$. Put $L(X)=\tr(X)/r$, $W_{AB}=L(P_AP_B)$ and $T_{ABC}=\operatorname{Re}L(P_AP_BP_C)$. Then
$W_{AB}\ge0$, $W_{AA}=1$, $W_{AB}=0$ on adjacent pairs, $T$ is symmetric, $T_{AAB}=W_{AB}$, and $T_{ABC}=0$ as
soon as one of the three pairs is adjacent. For a positive semidefinite operator $g$ and $X_0=1$, $X_B=P_B$
the matrix $[L(X_i^*gX_j)]$ is positive semidefinite. We use $g=1$ (\emph{two-point block}: $W-J/d\succeq0$),
$g=P_A$ (\emph{vertex block}: the matrix $[T_{ABC}]_{B,C}$ over the triples not adjacent to $A$), and
$g=1-\sum_{A\in Q}P_A$ for a clique $Q$ (\emph{clique block}; this $g$ is a projector). All conditions are preserved by averaging over
the symmetric group; after averaging $W$ takes a value $w_1$ on pairs sharing two elements and $w_3$ on
disjoint pairs, and $T$ takes four values $T_s$ (three triples through a common pair), $T_k$ (three triples in
a $4$-set), $T_p$ (two triples sharing a pair, the third disjoint from both), $T_d$ (three pairwise disjoint
triples).

Let $m=t-3$, $V=\binom t3$, $N_3=\binom m3$, and for $0\le w\le\tfrac13$
\[ \kappa_m=\frac{4m^4(m-1)}{9(m-3)},\qquad
F_m(w)=1+3mw+\tfrac12\Bigl[\tfrac m3+\sqrt{\tfrac{m^2}9+\kappa_m\,w(1-3w)}\,\Bigr]. \]

\begin{theorem}\label{thm:A}
Let $t\ge7$. Every projective representation of $\Gt$, over $\R$ or over $\C$ with the Hermitian form, has cost
at least $V/\max_{0\le w\le1/3}F_m(w)\ge V/\Phi(m)$, where
$\Phi(m)=1+\tfrac{2m}3+\sqrt{\bigl(\tfrac{3m^2}{\kappa_m}+\tfrac14\bigr)\bigl(\tfrac{m^2}9+\tfrac{\kappa_m}{12}\bigr)}$, and
$V/\Phi(m)=\bigl(\sqrt3-(12-2\sqrt3)/m+O(m^{-2})\bigr)\,t$.
\end{theorem}

\begin{proof}
A triple shares two elements with $3m$ others and is disjoint from $N_3$ others. We use five consequences of
the blocks.
(1) The two-point block tested on the constants: $d\ge V/\sigma$, $\sigma=1+3mw_1+N_3w_3$.
(2) Bessel's inequality for the clique $\{abp,acq,bcr,pqr\}$ tested against $abc$: $3w_1+w_3\le1$.
(3) The vertex block of a triple $A$, compressed to the span of $e_A$, of the indicator of the $3m$ triples
sharing two elements with $A$, and of the indicator of the $N_3$ triples disjoint from $A$, is the positive
semidefinite matrix
\[\begin{pmatrix}1&3mw_1&N_3w_3\\ 3mw_1&3m\bigl(w_1+(m-1)T_s+2T_k\bigr)&3m\tbinom{m-1}3T_p\\
N_3w_3&3m\tbinom{m-1}3T_p&N_3\bigl(w_3+3(m-3)T_p+\tbinom{m-3}3T_d\bigr)\end{pmatrix},\]
whose Schur complements give
(S1) $3m(w_1+(m-1)T_s+2T_k)\ge9m^2w_1^2$ and (S2) $w_3+3(m-3)T_p+\binom{m-3}3T_d\ge N_3w_3^2$.
(4) For a function $\varphi$ on the $m$ points outside $A$ with $\sum\varphi=0$, $\sum\varphi^2=1$, testing the
vertex block on $(e,x)\mapsto\varphi(x)/\sqrt3$ (triples $e\cup\{x\}$, $e$ a pair of $A$) and on
$C\mapsto\sum_{x\in C}\varphi(x)/\sqrt{\nu}$, $\nu=\binom{m-2}2$, gives a positive semidefinite $2\times2$ matrix
with diagonal entries $\mathcal A=w_1-T_s+2T_k$ and
$\mathcal B=w_3+(2m-9)T_p-\tfrac{(m-4)(m-5)}2T_d$ and off-diagonal entry $-\sqrt{3\nu}\,T_p$; hence
(D2) $\mathcal A\mathcal B\ge3\binom{m-2}2T_p^2$.
(5) The component of the vertex block with sum zero over the pairs of $A$ and over the outside points:
$w_1-T_s-T_k\ge0$.

By (5), $T_k\le w_1-T_s$; with (S1), $(m-3)T_s\ge3mw_1^2-3w_1$, hence
$\mathcal A\le3(w_1-T_s)\le3mw_1(1-3w_1)/(m-3)$. By (D2),
$\tfrac{(m-4)(m-5)}2T_d\le w_3+(2m-9)T_p-3\binom{m-2}2T_p^2/\mathcal A$ (if $\mathcal A=0$ then $T_p=0$ by
(D2), and the last term is read as $0$); multiplying by $(m-3)/3$, using
$\binom{m-3}3=\tfrac{m-3}3\cdot\tfrac{(m-4)(m-5)}2$ and comparing with (S2),
\[N_3w_3^2\le\tfrac m3w_3+\tfrac{2m(m-3)}3T_p-\tfrac{(m-3)\binom{m-2}2}{\mathcal A}T_p^2
\le\tfrac m3w_3+\frac{2m^2\mathcal A}{9(m-2)},\]
the last step being the maximum of a quadratic polynomial in $T_p$. With $y=N_3w_3$ this is
$y^2-\tfrac m3y-\tfrac{2N_3m^2\mathcal A}{9(m-2)}\le0$, so
$y\le\tfrac12\bigl[\tfrac m3+\sqrt{m^2/9+8N_3m^2\mathcal A/(9(m-2))}\bigr]$. The right-hand side increases
with $\mathcal A$; inserting the bound for
$\mathcal A$ and $8N_3m^3/(3(m-2)(m-3))=\kappa_m$ gives $\sigma\le F_m(w_1)$, and $0\le w_1\le\tfrac13$ by (2).
For the maximum put $\eta=w-\tfrac16$, so that $w(1-3w)=\tfrac1{12}-3\eta^2$; by the Cauchy--Schwarz inequality
$3m\eta+\tfrac12\sqrt{m^2/9+\kappa_m/12-3\kappa_m\eta^2}\le\sqrt{(3m^2/\kappa_m+\tfrac14)(m^2/9+\kappa_m/12)}$.
\end{proof}

\begin{theorem}\label{thm:tplus2}
$d_\R(t)\ge t+2$ and $d_{\mathrm H}(t)\ge t+2$ for every $t\ge11$. In particular the real dimension exceeds the
dimension $t+1$ of the complex bilinear representation for every $t\ge11$.
\end{theorem}
\begin{proof}
For $t=11,12$ this is Theorem~\ref{thm:cert}. Let $t\ge13$ and $m=t-3=10+s$, $s\ge0$. By
Theorem~\ref{thm:A} it suffices that $V/\Phi(m)>t+1$, that is $\mathcal L>\sqrt{\mathcal Q}$ with
$\mathcal L=V/(t+1)-1-2m/3$ and $\mathcal Q=(3m^2/\kappa_m+\tfrac14)(m^2/9+\kappa_m/12)$. One computes
$\mathcal L=(s^3+32s^2+329s+1072)/(6(s+14))$ and
$\mathcal L^2-\mathcal Q=\mathcal P(s)/\mathcal D(s)$ with
\begin{align*}
\mathcal P(s)&=s^8+77s^7+2528s^6+45938s^5+500519s^4+3293219s^3\\
&\quad+12363712s^2+22285350s+10306926,\\
\mathcal D(s)&=54s^4+2376s^3+38178s^2+264600s+666792.
\end{align*}
All coefficients are positive. So the cost exceeds $t+1$, and a dimension is an integer.
\end{proof}

\begin{theorem}\label{thm:cert}
For $t=10,11,\dots,16$ every projective representation of $\Gt$ has cost at least
$10.133$, $13.634$, $15.379$, $16.739$, $18.428$, $20.316$, $22.275$. Hence $d_\R(t)$ and $d_{\mathrm H}(t)$ are
at least $11,14,16,17,19,21,23$.
\end{theorem}
\begin{proof}
Each bound is the sum, with positive rational multipliers, of seven inequalities $u^{\mathsf T}M(x)u\ge0$,
where $M$ is the two-point block, the vertex block, or the clique block of a Fano plane or of a maximal star
of triples through a point, evaluated at $x=(d,w_1,w_3,T_s,T_k,T_p,T_d)$, and $u$ is an integer vector; the
weighted sum is the inequality $d-\beta\ge0$. The vectors and multipliers are in the files
\texttt{three\_point\_cert\_t\{10,\dots,16\}\_proj.json}; verifying a certificate needs integer arithmetic
only.
\end{proof}

\begin{proposition}\label{prop:proj}
$G_{10}$ has a projective representation of rank $2$ in dimension $25$ (cost $25/2$). Hence no argument that
is valid for projectors of every rank proves more than $12.5$ at $t=10$, where the best representation by
vectors that we know has dimension $15$.
\end{proposition}
\begin{proof}
For a triple $A$ inside the first nine points let $\rho_A=1_A-\tfrac13\mathbf 1\in\R^9$; then
$\ip{\rho_A}{\rho_B}=|A\cap B|-1$, and the $\rho_A$ span a space $E_0$ of dimension $8$. Give $A$ the plane
$\rho_A\otimes\R^2$ in $E_0\otimes\R^2$ (dimension $16$), and give a triple $\{a,b,10\}$ the coordinate plane of
$e_a,e_b$ in a further $\R^9$. Two triples through the point $10$ are adjacent exactly when their pairs are
disjoint, and planes in different summands are orthogonal.
\end{proof}

By Corollary~\ref{cor:proj} the projective cost is nevertheless of order $t^2$; the moment method of this
section simply does not see it.

\section{Formal verification}\label{sec:lean}
% Section "Formal verification" of PAPER_TR26_151_PROBLEM1.tex (same header and claim boundary as the main
% file). Source of every statement in this section: lean_theorem_l/BUILD_RECORD.md, axioms_output.txt,
% axioms_output_layer_ranks.txt, verify_by_main_session.txt and verify_layer_ranks_by_main_session.txt
% (independent re-runs), referee_g10/ (audit of the statements of TheoremL.lean). Source of the paragraph
% "Clean-machine build" (added on 27 November 2026): public_release_tr26_151/lean/records/ci_run_1.md.
Theorem~\ref{thm:L} and Corollaries~\ref{cor:triples} and~\ref{cor:proj} have been formalised in Lean~4 (version
4.30.0-rc2) with Mathlib (tag \texttt{v4.30.0-rc2}), in two files of the folder \texttt{lean\_theorem\_l}:
\texttt{TheoremL.lean} ($1065$ lines, SHA-256 \texttt{3019bee1\,\dots\,6b3cef92}) and
\texttt{TheoremL\_LayerRanks.lean} ($1354$ lines, SHA-256 \texttt{a1b16696\,\dots\,21adeefd}), whose first
$1065$ lines are the first file, byte for byte. The full hashes are in the build record. Neither file contains
\texttt{sorry}, an added axiom or \texttt{native\_decide}. The main declarations are:
\begin{center}\small
\begin{tabular}{@{}p{0.39\textwidth}p{0.58\textwidth}@{}}
\toprule
declaration & statement\\
\midrule
\texttt{core\_bound} & orthonormal $r$-frames $u(p,\cdot)$ in a real space $E$, orthogonal on adjacent
points: $r\,n^2\le5\dim E$\\
\texttt{core\_layer\_ranks}$^{\ast}$ & the same hypotheses: $3\,r\,n^3\le5\sum_\Lambda\rank(\Lambda)$
(Theorem~\ref{thm:L} over $\R$)\\
\texttt{theorem\_L\_real} & unit vectors in $\R^d$: $n^2\le5d$ \ (\texttt{\_nine}: $n^2\le9d$)\\
\texttt{theorem\_L\_complex\_hermitian} & unit vectors in $\C^d$, Hermitian form: $n^2\le5d$\\
\texttt{theorem\_L\_projective} & subspaces of dimension $r$ of a real space: $r\,n^2\le5\dim E$\\
\texttt{theorem\_L\_projective\_complex} & the same over $\C$\\
\texttt{theorem\_L\_projectors} & self-adjoint operators on $\R^D$ with ranges of dimension $r$ and
$P_pP_q=0$ on adjacent points: $r\,n^2\le5D$\\
\texttt{corollary\_L\_real} & non-zero vectors in $\R^d$ indexed by the $3$-subsets of $[t]$, orthogonal
when two subsets share one element: $\floor{t/3}^2\le5d$ \ (\texttt{\_nine}: $\le9d$)\\
\texttt{corollary\_L\_complex\_hermitian}$^{\ast}$ & the same in $\C^d$ with the Hermitian form\\
\texttt{corollary\_L\_projective\_complex}$^{\ast}$ & complex subspaces of dimension $r$ indexed by the
$3$-subsets, orthogonal when two subsets share one element: $r\floor{t/3}^2\le5\dim_\C E$
(Corollary~\ref{cor:proj} for ranges)\\
\texttt{upper\_bound\_triples}$^{\ast}$ & a representation of $\Gt$ by non-zero vectors of $\R^{t^2}$ exists\\
\bottomrule
\end{tabular}
\end{center}
($^{\ast}$: in the second file.) The statement about the quantity of Problem~1 of~\cite{HI} reads, with the
Unicode symbols of Lean transliterated:
\begin{verbatim}
theorem corollary_L_real {t d : Nat}
    (v : {A : Finset (Fin t) // A.card = 3} -> EuclideanSpace Real (Fin d))
    (h_ne : forall A, v A != 0)
    (h_orth : forall A B, (A.1 inter B.1).card = 1 ->
                inner Real (v A) (v B) = 0) :
    (t / 3) ^ 2 <= 5 * d
\end{verbatim}
Here \texttt{t / 3} is the division of natural numbers. Together with \texttt{upper\_bound\_triples} (a
representation by non-zero vectors indexed by $\texttt{Fin}\ t\times\texttt{Fin}\ t$, that is, in dimension
$t^2$) it gives $\floor{t/3}^2/5\le d_\R(t)\le t^2$; the minimum $d_\R(t)$ itself is not defined in the Lean
files.

\paragraph{Checks.} Each file was compiled by the Lean kernel twice, by the process that wrote it and in an
independent re-run, with exit code $0$ and no warning. For every declaration of the table,
\texttt{\#print axioms} reports exactly \texttt{propext}, \texttt{Classical.choice} and \texttt{Quot.sound}.
As a negative control, a file that states \texttt{theorem\_L\_real} and proves it by \texttt{sorry} is
reported as depending on \texttt{sorryAx}. A further theorem shows that the hypotheses of
\texttt{theorem\_L\_real} can be satisfied for every $n$, so that statement is not vacuous.
% Result of the isolated audit of the Lean statements (referee_g10/REFEREE_REPLY.md, 26 November 2026).
The formal statements were then compared with the definitions of~\cite{HI} (its pages 3--4 and~14) by an
isolated reader who did not run Lean. The reader found \texttt{corollary\_L\_real} to state exactly the lower
bound asked for in Problem~1 (vertices: the $3$-subsets; non-zero vectors of $\R^d$ with the standard inner
product; orthogonality required exactly for two subsets that share one element; no further hypothesis), and
found no construct in the file that changes the meaning of a statement. The reader's caveats concern
coverage only and are listed in the next paragraph. A second isolated reader audited the declarations of the
second file in the same way, with the same conclusion.


\paragraph{What is not formalised.} The formal proof follows Section~\ref{sec:proof} in the frame form
($u(p,k)$ in place of $U_p$) and over $\R$; the complex statements are derived inside Lean from the real ones
(a complex subspace of dimension $r$ is a real subspace of dimension $2r$), not by repeating the proof over
$\C$. The following are \emph{not} formalised: the definitions of $\xi_f$ and $\chi_q$ and the inequality
$\xi_f\le\chi_q$ quoted from~\cite{MR}; the passage from complex projectors with $P_AP_B=0$ to orthogonal
ranges (formalised over $\R$ only); Remark~\ref{rem:k}; the statements of Section~\ref{sec:mink}; and all of
Section~\ref{sec:linear} (the certificates of Theorem~\ref{thm:cert} are checked by exact-arithmetic scripts,
not by Lean). The declarations marked $^{\ast}$ were added after the first audit; a second isolated
reader then compared them with the mathematics (\texttt{referee\_g11}) and found them aligned, with caveats
about coverage that are all listed in this paragraph.

\paragraph{Clean-machine build (added on 27 November 2026).} The checks above were run against an existing
Mathlib build, on the machine on which the files were written. The build kit for a clean machine
(\texttt{lakefile.toml}, \texttt{lean-toolchain} and the workflow \texttt{lean\_verify}) was then executed on a
fresh GitHub-hosted runner, by GitHub Actions (run~1 of the workflow, 27~November 2026), and every step
succeeded. The build of the two files printed $27$ lines of \texttt{\#print axioms} output ($10$ from the
first file, $17$ from the second), each naming exactly \texttt{propext}, \texttt{Classical.choice} and
\texttt{Quot.sound}; the planted \texttt{sorry} of the negative control was reported as \texttt{sorryAx}; and
the exact-arithmetic verifier accepted all seven certificates of Theorem~\ref{thm:cert}. The record of the run
is the file \texttt{lean/records/ci\_run\_1.md} of the repository
\url{https://github.com/aqida-ai/tr26-151-problem1}. That run used the same version of Lean and the same
Mathlib as the first checks, and the Mathlib binaries were downloaded from the Mathlib cache, not rebuilt. No
proof checker other than the Lean kernel was used.


\section{Related work, remarks and open questions}\label{sec:lit}
% Section "Related work, remarks and open questions" of PAPER_TR26_151_PROBLEM1.tex (same header and claim
% boundary as the main file). Status of the citations: [HI] read in full; [Ha] statement read; [MR], [CMNSW],
% [GdLL]: statements reported by a literature agent that read the arXiv versions (QUANTUM_PARAMETERS_CHECK_
% 20261126.md); [LCLL], [HT], [GH]: reported by a search agent, abstracts or snippets only
% (NOVELTY_SEARCH_20261126.md). None of the last three is used in a proof.

\paragraph{The problem and the earlier bounds.}
Problem~1 is stated in Section~5 of~\cite{HI}; the report page showed no revision and no comment on
26~November 2026. The lower bound of~\cite{HI} for real representations of Cayley-type graphs is a linear
programme which, by their Lemma~22, equals the Lov\'asz number. Bounds beyond the Lov\'asz number for the
dimension of orthogonal representations are known from topological arguments~\cite{Ha} (for Kneser-type
graphs; see also~\cite{GH}). Any argument that does not use the positivity of the form is limited to $t+1$ for
$\Gt$, by the complex bilinear representation; Section~\ref{sec:mink} locates the single place where the
proof of Theorem~\ref{thm:L} uses it.

\paragraph{Projective rank and quantum chromatic number.}
Following Man\v{c}inska and Roberson~\cite[Section~6]{MR}, a $d/r$-representation of a graph assigns to each
vertex a projector of rank $r$ in $M_d(\C)$ so that adjacent vertices get orthogonal projectors; the projective
rank $\xi_f$ is the infimum of $d/r$. They prove $\bar\vartheta\le\xi_f$ \cite[Theorem~6.4]{MR} and
$\xi_f\le\chi_q$ \cite[Lemma~6.2]{MR} (numbering of arXiv version~2), where $\chi_q$ is the quantum chromatic number of~\cite{CMNSW}; see also
\cite[Section~3]{GdLL}. Both parameters are monotone under taking subgraphs. Corollary~\ref{cor:proj} is
Theorem~\ref{thm:L} over $\C$ read with these definitions. Together with $\chi_q\le\chi\le\binom t2$ it gives
$\chi_q(J(t,3,1))=\Theta(t^2)$ and $\xi_f(J(t,3,1))=\Theta(t^2)$, while $\bar\vartheta(J(t,3,1))$ is of order
$t$: for this family the projective rank is far from the Lov\'asz bound. Exact quantum chromatic numbers of
generalised Johnson graphs in other parameter ranges are the subject of~\cite{LCLL}; minimum-rank parameters
of complements of threshold Kneser graphs are computed in~\cite{HT}.

\paragraph{Novelty.}
A search of the literature made on 26~November 2026 (queries on the cube graph under several names, on
orthogonal, minimum semidefinite and projective rank of Hamming and Johnson-type graphs, and on citations
of~\cite{HI}) found no earlier bound of order $n^2$ for $\Hn$ or of order $t^2$ for $\Gt$. A search that
finds nothing does not prove novelty.

\paragraph{Remarks and open questions.}
\begin{enumerate}
\item \emph{Constants.} For the cube graph, $n^2/5\le\dim\le n^2$. Its independence number is
$\max(n,2n-2,3n-6)$ (three pieces of lines in the three directions), so colourings need at least
$n^3/(3n-6)$ colours when $n\ge4$; at $n=4$ the chromatic number is $11$. Representations that are invariant under a product
$\Gamma_1\times\Gamma_2\times\Gamma_3$ of abelian groups acting regularly on the three coordinates need at least
$2n^2/3$ dimensions; over $\C$ this is attained when $3$ divides $n$ (over $\R$ we checked $n=3,6,9$).
Invariant representations are not optimal: at $n=4$ they need $14$ dimensions, and a colouring needs $11$. For $\Gt$ the truth
lies between $\floor{t/3}^2/5$ and $\binom t2$. What is the right constant?
\item \emph{Small cases.} $d_\R(9)=8$, and $11\le d_\R(10)\le15$. Every representation of $G_{10}$ whose
restriction to some nine points is the $E_8$ configuration has dimension at least $15$: the vector of
$\{a,b,10\}$ is orthogonal to the vectors of the triples of the nine points that meet $\{a,b\}$ once, and
these span the whole $8$-dimensional space (a rank computation on integers); so the vectors of the triples
through the point $10$ lie in the orthogonal complement of that space, and they need $7$ dimensions by
Haviv's theorem for the Kneser graph of pairs \cite[Corollary~3.3]{Ha}. We do not know $d_\R(10)$.
\item \emph{Other problems of~\cite{HI}.} We do not see how to apply the argument to their Problems~2 and~3.
For Problem~2 (vectors of $\{\pm1\}^n$, orthogonal at Hamming distance divisible by four) we tried the version
of the argument in which lines and layers are cosets of subgroups of $\mathbb F_2^n$; there the hypotheses
force the sum of two of the subgroups to be a doubly even code, and the resulting bound stays below the
clique bound $2^{n/2}$ (this computation was not refereed). For
Problem~3 (maps of the sphere that send pairs with inner product $\delta$ to orthogonal vectors) the argument
applies to any copy of a cube graph inside the graph of the sphere and bounds the dimension from below by one
fifth of the square of its side; we found no configuration for which this improves on the exponential bounds of~\cite{HI}.
\item \emph{The moment method.} The bounds of Section~\ref{sec:linear} are obtained from moments of order at
most three and hold for projectors; by Theorem~\ref{thm:L} the projective cost is of order $t^2$, so that
hierarchy is far from the truth at low order. Does a fixed level of it give a superlinear bound?
\end{enumerate}


\paragraph{Statement on the production of this manuscript.}
The mathematics, the computations and the text were produced with an AI system (Claude) under the direction
of the author, in a run that kept a written control file with a gate before each attempt. Intermediate
versions of every result were read by isolated AI referees with their own code. No human mathematician had
checked the manuscript when it was written.

\begin{thebibliography}{99}
\bibitem{HI} P.~Hrube\v{s} and S.~Iyer, \emph{Sparse polynomials and orthogonal representations of
combinatorial graphs}, Electronic Colloquium on Computational Complexity, Report TR26-151, 2026.
\bibitem{H24} P.~Hrube\v{s}, \emph{A subquadratic upper bound on sum-of-squares composition formulas},
Computational Complexity Conference (CCC), 2024.
\bibitem{Lo} L.~Lov\'asz, \emph{On the Shannon capacity of a graph}, IEEE Trans.\ Inform.\ Theory 25 (1979),
1--7.
% Further bibliography entries of PAPER_TR26_151_PROBLEM1.tex. Reading status: see the comment at the top of
% PAPER_TR26_151_PROBLEM1_lit.tex.
\bibitem{Ha} I.~Haviv, \emph{Topological bounds on the dimension of orthogonal representations of graphs},
European J.\ Combin.\ 81 (2019), 84--97; arXiv:1811.11488.
\bibitem{GH} A.~Golovnev and I.~Haviv, \emph{The (generalized) orthogonality dimension of (generalized) Kneser
graphs: bounds and applications}, Theory Comput.\ 18 (2022); arXiv:2002.08580.
\bibitem{MR} L.~Man\v{c}inska and D.~E.~Roberson, \emph{Quantum homomorphisms}, J.\ Combin.\ Theory Ser.~B
118 (2016); arXiv:1212.1724v2 (the statement numbers cited here, Lemma~6.2 and Theorem~6.4, are those of
this arXiv version).
\bibitem{CMNSW} P.~J.~Cameron, A.~Montanaro, M.~W.~Newman, S.~Severini and A.~Winter, \emph{On the quantum
chromatic number of a graph}, Electron.\ J.\ Combin.\ 14 (2007), R81.
\bibitem{GdLL} S.~Gribling, D.~de~Laat and M.~Laurent, \emph{Bounds on entanglement dimensions and quantum
graph parameters via noncommutative polynomial optimization}, Math.\ Program.\ 170 (2018); arXiv:1708.09696.
\bibitem{LCLL} G.~Luo, X.~Cao, S.~Li and Y.~Li, \emph{On the exact quantum chromatic number of generalized
Johnson graphs}, arXiv:2609.04677, 2026.
\bibitem{HT} T.~Hu and Q.~Tang, \emph{Minimum-rank parameters of complements of threshold Kneser graphs},
arXiv:2607.06480, 2026.

\end{thebibliography}
\end{document}
